% !TEX program = xelatex
% À compiler sur Overleaf avec le compilateur XeLaTeX
% (Menu -> Compiler -> XeLaTeX) pour avoir la police Verdana.

\documentclass[12pt,a4paper]{article}

\usepackage{fontspec}
% Verdana si disponible, sinon une police sans-serif proche
\IfFontExistsTF{Verdana}{\setmainfont{Verdana}}{\setmainfont{DejaVu Sans}}

\usepackage[french]{babel}
\usepackage[margin=1.4cm]{geometry}
\usepackage{amsmath}
\usepackage{setspace}
\usepackage{enumitem}
\usepackage{array}

% Interligne 1,5
\onehalfspacing

% Espacement réduit autour des égalités successives
\setlength{\abovedisplayskip}{1pt}
\setlength{\belowdisplayskip}{1pt}
\setlength{\abovedisplayshortskip}{0pt}
\setlength{\belowdisplayshortskip}{0pt}

% ----- Pas de numérotation des pages -----
\pagestyle{empty}

\setlength{\parindent}{0pt}

% ----- Commandes personnalisées -----
\newcounter{exercice}
\newcommand{\exercice}{%
  \refstepcounter{exercice}%
  \vspace{0.4em}
  \noindent\underline{\textbf{Exercice \arabic{exercice}}}\par
  \vspace{0em}
}

\setlist[enumerate,1]{label=\alph*),itemsep=0pt,topsep=1pt}
\setlist[itemize,1]{leftmargin=1.6em,itemsep=0pt,topsep=1pt}

\begin{document}

% =====================================================================
% VERSION 1
% =====================================================================
\setcounter{exercice}{0}

% ----- En-tête -----
\begin{center}
  \underline{\Large\textbf{Correction du contrôle séquence 13 version 1}}\\[0.3em]
 
\end{center}

% =====================================================================
\vspace{3cm}\exercice
\begin{enumerate}
  \item \textbf{Oui.} $36  = 6 \times 6$ 
  \newline 6 est donc un diviseur de 36.
  \item \textbf{Non.} $148 \div 7 \approx 21{,}14$
  \newline le quotient n'est pas un entier,
        donc 148 n'est pas un multiple de 7.
  \item \textbf{Non.}  $10 < 20$, donc 10 n'est pas un multiple de 20.
  \item \textbf{Non.} $38 \div 4 = 9{,}5$ 
  \newline le quotient n'est pas un entier,
        donc 4 n'est pas un diviseur de 38.
\end{enumerate}

% =====================================================================
\vspace{1cm}\exercice
\begin{enumerate}
  \item \textbf{2 et 3} (acceptés aussi : 5 ; 7 ; 11 ; \dots).
        Un nombre premier a exactement deux diviseurs : 1 et lui-même.
  \item \textbf{4 et 6} (acceptés aussi : 1 ; 8 ; 9 ; \dots).
      
  \item Dans 12 ; 16 ; 11 ; 25 ; 49 ; 56 ; 83 ; 93, les deux nombres premiers sont
        \textbf{11} et \textbf{83}.
       
\end{enumerate}

% =====================================================================
\vspace{1cm}\exercice
\textbf{a) } $6 = 2 \times 3$, 
donc les diviseurs de 6 sont $\boxed{1\ ;\ 6\ ;\ 2\ ;\ 3\ }$.

\textbf{b) } $20 = 2 \times 2 \times 5$,
donc les diviseurs de 20 sont
$\boxed{1\ ;\ 20\ ;\ 2\ ;\ 5\ ;\ 4\ ;\ 10\ }$.

% =====================================================================
\vspace{1cm}\exercice


\noindent\begin{minipage}[t]{0.48\textwidth}
\textbf{a)}
\begin{align*}
14 &= 2 \times 7
\end{align*}

\textbf{b)}
\begin{align*}
30 &= 2 \times 15\\
   &= 2 \times 3 \times 5
\end{align*}
\end{minipage}\hfill
\begin{minipage}[t]{0.48\textwidth}
\textbf{c)}
\begin{align*}
60 &= 2 \times 30\\
   &= 2 \times 2 \times 15\\
   &= 2 \times 2 \times 3 \times 5\\
\end{align*}

\textbf{d)}
\begin{align*}
105 &= 3 \times 35\\
    &= 3 \times 5 \times 7
\end{align*}
\end{minipage}

% =====================================================================
\exercice
\noindent\begin{minipage}[t]{0.31\textwidth}
\textbf{a)} 
\begin{align*}
\dfrac{3}{12} &= \dfrac{3 \times 1}{3 \times 4}\\
              &= \dfrac{1}{4}
\end{align*}
\end{minipage}\hfill
\begin{minipage}[t]{0.31\textwidth}
\textbf{b)} 
\begin{align*}
\dfrac{16}{24} &= \dfrac{8 \times 2}{8 \times 3}\\
               &= \dfrac{2}{3}
\end{align*}
\end{minipage}\hfill
\begin{minipage}[t]{0.31\textwidth}
\textbf{c)} 
\begin{align*}
\dfrac{60}{105} &= \dfrac{15 \times 4}{15 \times 7}\\
                &= \dfrac{4}{7}
\end{align*}
\end{minipage}

% =====================================================================
\vspace{1cm}\exercice
Le nombre de paquets doit diviser à la fois 105 et 60 : c'est un diviseur commun à
105 et 60. Pour en faire un maximum, on cherche le plus grand diviseur commun à 105 et 60.

Avec les décompositions en facteurs premiers obtenues à l'exercice 4 :
$105 = 3 \times 5 \times 7$ et $60 = 2 \times 2 \times 3 \times 5$.
Le plus grand diviseur commun à 105 et 60 est le produit des facteurs premiers communs :
3 \times 5
                      = 15

Hulk et Batman peuvent donc faire au maximum \textbf{15 paquets},
chacun contenant  7 cookies et 4 tartelettes.

% =====================================================================
\vspace{1cm}\exercice
On utilise les critères de divisibilité :
\newline \textbf{par 2} : le chiffre des unités est 0, 2, 4, 6 ou 8 ;
\newline \textbf{par 3} : la somme des chiffres est divisible par 3
(sommes des chiffres : $19$ pour 254\,350 ; $9$ pour 111\,111\,111 ; $22$ pour 811\,255) ;
\newline \textbf{par 5} : le chiffre des unités est 0 ou 5 ;
\newline \textbf{par 10} : le chiffre des unités est 0.

\renewcommand{\arraystretch}{1.4}
\setlength{\tabcolsep}{2pt}
\begin{center}
\small
\begin{tabular}{|c|c|c|c|c|}
\hline
 & \textbf{Divisible par 2} & \textbf{Divisible par 3} & \textbf{Divisible par 5} & \textbf{Divisible par 10} \\
\hline
254\,350 & Oui & Non & Oui & Oui \\
\hline
111\,111\,111 & Non & Oui & Non & Non \\
\hline
811\,255 & Non & Non & Oui & Non \\
\hline
\end{tabular}
\end{center}

\newpage

% =====================================================================
% VERSION 2
% =====================================================================
\setcounter{exercice}{0}

% ----- En-tête -----
\begin{center}
  \underline{\Large\textbf{Correction du contrôle séquence 13 version 2}}\\[0.3em]
 
\end{center}

% =====================================================================
\vspace{3cm}\exercice
\begin{enumerate}
  \item \textbf{Oui.} $42  = 7 \times 6$ 
  \newline 7 est donc un diviseur de 42.
  \item \textbf{Non.} $165 \div 8 \approx 20{,}63$
  \newline le quotient n'est pas un entier,
        donc 165 n'est pas un multiple de 8.
  \item \textbf{Non.}  $15 < 45$, donc 15 n'est pas un multiple de 45.
  \item \textbf{Non.} $44 \div 6 \approx 7{,}33$
  \newline le quotient n'est pas un entier,
        donc 6 n'est pas un diviseur de 44.
\end{enumerate}

% =====================================================================
\vspace{1cm}\exercice
\begin{enumerate}
  \item \textbf{2 et 3} (acceptés aussi : 5 ; 7 ; 11 ; \dots).
        Un nombre premier a exactement deux diviseurs : 1 et lui-même.
  \item \textbf{4 et 6} (acceptés aussi : 1 ; 8 ; 9 ; \dots).
      
  \item Dans 14 ; 15 ; 7 ; 27 ; 51 ; 62 ; 71 ; 91, les deux nombres premiers sont
        \textbf{7} et \textbf{71}.
       
\end{enumerate}

% =====================================================================
\vspace{1cm}\exercice
\textbf{a) } $8 = 2 \times 2 \times 2$, 
donc les diviseurs de 8 sont $\boxed{1\ ;\ 8\ ;\ 2\ ;\ 4\ }$.

\textbf{b) } $12 = 2 \times 2 \times 3$,
donc les diviseurs de 12 sont
$\boxed{1\ ;\ 12\ ;\ 2\ ;\ 3\ ;\ 4\ ;\ 6\ }$.

% =====================================================================
\vspace{1cm}\exercice


\noindent\begin{minipage}[t]{0.48\textwidth}
\textbf{a)}
\begin{align*}
15 &= 3 \times 5
\end{align*}

\textbf{b)}
\begin{align*}
42 &= 2 \times 21\\
   &= 2 \times 3 \times 7
\end{align*}
\end{minipage}\hfill
\begin{minipage}[t]{0.48\textwidth}
\textbf{c)}
\begin{align*}
84 &= 2 \times 42\\
   &= 2 \times 2 \times 21\\
   &= 2 \times 2 \times 3 \times 7\\
\end{align*}

\textbf{d)}
\begin{align*}
165 &= 3 \times 55\\
    &= 3 \times 5 \times 11
\end{align*}
\end{minipage}

% =====================================================================
\exercice
\noindent\begin{minipage}[t]{0.31\textwidth}
\textbf{a)} 
\begin{align*}
\dfrac{2}{10} &= \dfrac{2 \times 1}{2 \times 5}\\
              &= \dfrac{1}{5}
\end{align*}
\end{minipage}\hfill
\begin{minipage}[t]{0.31\textwidth}
\textbf{b)} 
\begin{align*}
\dfrac{14}{21} &= \dfrac{7 \times 2}{7 \times 3}\\
               &= \dfrac{2}{3}
\end{align*}
\end{minipage}\hfill
\begin{minipage}[t]{0.31\textwidth}
\textbf{c)} 
\begin{align*}
\dfrac{45}{165} &= \dfrac{15 \times 3}{15 \times 11}\\
                &= \dfrac{3}{11}
\end{align*}
\end{minipage}

% =====================================================================
\vspace{1cm}\exercice
Le nombre de paquets doit diviser à la fois 120 et 75 : c'est un diviseur commun à
120 et 75. Pour en faire un maximum, on cherche le plus grand diviseur commun à 120 et 75.

Avec les décompositions en facteurs premiers :
$120 = 2 \times 2 \times 2 \times 3 \times 5$ et $75 = 3 \times 5 \times 5$.
Le plus grand diviseur commun à 120 et 75 est le produit des facteurs premiers communs :
3 \times 5
                      = 15

Hulk et Batman peuvent donc faire au maximum \textbf{15 paquets},
chacun contenant  8 cookies et 5 tartelettes.

% =====================================================================
\vspace{1cm}\exercice
On utilise les critères de divisibilité :
\newline \textbf{par 2} : le chiffre des unités est 0, 2, 4, 6 ou 8 ;
\newline \textbf{par 3} : la somme des chiffres est divisible par 3
(sommes des chiffres : $25$ pour 372\,940 ; $18$ pour 222\,222\,222 ; $28$ pour 624\,835) ;
\newline \textbf{par 5} : le chiffre des unités est 0 ou 5 ;
\newline \textbf{par 10} : le chiffre des unités est 0.

\renewcommand{\arraystretch}{1.4}
\setlength{\tabcolsep}{2pt}
\begin{center}
\small
\begin{tabular}{|c|c|c|c|c|}
\hline
 & \textbf{Divisible par 2} & \textbf{Divisible par 3} & \textbf{Divisible par 5} & \textbf{Divisible par 10} \\
\hline
372\,940 & Oui & Non & Oui & Oui \\
\hline
222\,222\,222 & Oui & Oui & Non & Non \\
\hline
624\,835 & Non & Non & Oui & Non \\
\hline
\end{tabular}
\end{center}

\newpage

% =====================================================================
% VERSION 3
% =====================================================================
\setcounter{exercice}{0}

% ----- En-tête -----
\begin{center}
  \underline{\Large\textbf{Correction du contrôle séquence 13 version 3}}\\[0.3em]
 
\end{center}

% =====================================================================
\vspace{3cm}\exercice
\begin{enumerate}
  \item \textbf{Oui.} $56  = 8 \times 7$ 
  \newline 8 est donc un diviseur de 56.
  \item \textbf{Non.} $200 \div 9 \approx 22{,}22$
  \newline le quotient n'est pas un entier,
        donc 200 n'est pas un multiple de 9.
  \item \textbf{Non.}  $14 < 42$, donc 14 n'est pas un multiple de 42.
  \item \textbf{Non.} $50 \div 9 \approx 5{,}56$
  \newline le quotient n'est pas un entier,
        donc 9 n'est pas un diviseur de 50.
\end{enumerate}

% =====================================================================
\vspace{1cm}\exercice
\begin{enumerate}
  \item \textbf{2 et 3} (acceptés aussi : 5 ; 7 ; 11 ; \dots).
        Un nombre premier a exactement deux diviseurs : 1 et lui-même.
  \item \textbf{4 et 6} (acceptés aussi : 1 ; 8 ; 9 ; \dots).
      
  \item Dans 10 ; 21 ; 5 ; 33 ; 55 ; 68 ; 89 ; 87, les deux nombres premiers sont
        \textbf{5} et \textbf{89}.
       
\end{enumerate}

% =====================================================================
\vspace{1cm}\exercice
\textbf{a) } $10 = 2 \times 5$, 
donc les diviseurs de 10 sont $\boxed{1\ ;\ 10\ ;\ 2\ ;\ 5\ }$.

\textbf{b) } $50 = 2 \times 5 \times 5$,
donc les diviseurs de 50 sont
$\boxed{1\ ;\ 50\ ;\ 2\ ;\ 5\ ;\ 10\ ;\ 25\ }$.

% =====================================================================
\vspace{1cm}\exercice


\noindent\begin{minipage}[t]{0.48\textwidth}
\textbf{a)}
\begin{align*}
22 &= 2 \times 11
\end{align*}

\textbf{b)}
\begin{align*}
70 &= 2 \times 35\\
   &= 2 \times 5 \times 7
\end{align*}
\end{minipage}\hfill
\begin{minipage}[t]{0.48\textwidth}
\textbf{c)}
\begin{align*}
132 &= 2 \times 66\\
    &= 2 \times 2 \times 33\\
    &= 2 \times 2 \times 3 \times 11\\
\end{align*}

\textbf{d)}
\begin{align*}
195 &= 3 \times 65\\
    &= 3 \times 5 \times 13
\end{align*}
\end{minipage}

% =====================================================================
\exercice
\noindent\begin{minipage}[t]{0.31\textwidth}
\textbf{a)} 
\begin{align*}
\dfrac{4}{20} &= \dfrac{4 \times 1}{4 \times 5}\\
              &= \dfrac{1}{5}
\end{align*}
\end{minipage}\hfill
\begin{minipage}[t]{0.31\textwidth}
\textbf{b)} 
\begin{align*}
\dfrac{12}{18} &= \dfrac{6 \times 2}{6 \times 3}\\
               &= \dfrac{2}{3}
\end{align*}
\end{minipage}\hfill
\begin{minipage}[t]{0.31\textwidth}
\textbf{c)} 
\begin{align*}
\dfrac{75}{105} &= \dfrac{15 \times 5}{15 \times 7}\\
                &= \dfrac{5}{7}
\end{align*}
\end{minipage}

% =====================================================================
\vspace{1cm}\exercice
Le nombre de paquets doit diviser à la fois 90 et 105 : c'est un diviseur commun à
90 et 105. Pour en faire un maximum, on cherche le plus grand diviseur commun à 90 et 105.

Avec les décompositions en facteurs premiers :
$90 = 2 \times 3 \times 3 \times 5$ et $105 = 3 \times 5 \times 7$.
Le plus grand diviseur commun à 90 et 105 est le produit des facteurs premiers communs :
3 \times 5
                      = 15

Hulk et Batman peuvent donc faire au maximum \textbf{15 paquets},
chacun contenant  6 cookies et 7 tartelettes.

% =====================================================================
\vspace{1cm}\exercice
On utilise les critères de divisibilité :
\newline \textbf{par 2} : le chiffre des unités est 0, 2, 4, 6 ou 8 ;
\newline \textbf{par 3} : la somme des chiffres est divisible par 3
(sommes des chiffres : $23$ pour 483\,260 ; $27$ pour 333\,333\,333 ; $29$ pour 367\,445) ;
\newline \textbf{par 5} : le chiffre des unités est 0 ou 5 ;
\newline \textbf{par 10} : le chiffre des unités est 0.

\renewcommand{\arraystretch}{1.4}
\setlength{\tabcolsep}{2pt}
\begin{center}
\small
\begin{tabular}{|c|c|c|c|c|}
\hline
 & \textbf{Divisible par 2} & \textbf{Divisible par 3} & \textbf{Divisible par 5} & \textbf{Divisible par 10} \\
\hline
483\,260 & Oui & Non & Oui & Oui \\
\hline
333\,333\,333 & Non & Oui & Non & Non \\
\hline
367\,445 & Non & Non & Oui & Non \\
\hline
\end{tabular}
\end{center}

\newpage

% =====================================================================
% VERSION 4
% =====================================================================
\setcounter{exercice}{0}

% ----- En-tête -----
\begin{center}
  \underline{\Large\textbf{Correction du contrôle séquence 13 version 4}}\\[0.3em]
 
\end{center}

% =====================================================================
\vspace{3cm}\exercice
\begin{enumerate}
  \item \textbf{Oui.} $54  = 9 \times 6$ 
  \newline 9 est donc un diviseur de 54.
  \item \textbf{Non.} $126 \div 5 = 25{,}2$
  \newline le quotient n'est pas un entier,
        donc 126 n'est pas un multiple de 5.
  \item \textbf{Non.}  $12 < 36$, donc 12 n'est pas un multiple de 36.
  \item \textbf{Non.} $48 \div 5 = 9{,}6$ 
  \newline le quotient n'est pas un entier,
        donc 5 n'est pas un diviseur de 48.
\end{enumerate}

% =====================================================================
\vspace{1cm}\exercice
\begin{enumerate}
  \item \textbf{2 et 3} (acceptés aussi : 5 ; 7 ; 11 ; \dots).
        Un nombre premier a exactement deux diviseurs : 1 et lui-même.
  \item \textbf{4 et 6} (acceptés aussi : 1 ; 8 ; 9 ; \dots).
      
  \item Dans 12 ; 18 ; 3 ; 35 ; 57 ; 64 ; 97 ; 77, les deux nombres premiers sont
        \textbf{3} et \textbf{97}.
       
\end{enumerate}

% =====================================================================
\vspace{1cm}\exercice
\textbf{a) } $15 = 3 \times 5$, 
donc les diviseurs de 15 sont $\boxed{1\ ;\ 15\ ;\ 3\ ;\ 5\ }$.

\textbf{b) } $45 = 3 \times 3 \times 5$,
donc les diviseurs de 45 sont
$\boxed{1\ ;\ 45\ ;\ 3\ ;\ 5\ ;\ 9\ ;\ 15\ }$.

% =====================================================================
\vspace{1cm}\exercice


\noindent\begin{minipage}[t]{0.48\textwidth}
\textbf{a)}
\begin{align*}
21 &= 3 \times 7
\end{align*}

\textbf{b)}
\begin{align*}
20 &= 2 \times 10\\
   &= 2 \times 2 \times 5
\end{align*}
\end{minipage}\hfill
\begin{minipage}[t]{0.48\textwidth}
\textbf{c)}
\begin{align*}
156 &= 2 \times 78\\
    &= 2 \times 2 \times 39\\
    &= 2 \times 2 \times 3 \times 13\\
\end{align*}

\textbf{d)}
\begin{align*}
231 &= 3 \times 77\\
    &= 3 \times 7 \times 11
\end{align*}
\end{minipage}

% =====================================================================
\exercice
\noindent\begin{minipage}[t]{0.31\textwidth}
\textbf{a)} 
\begin{align*}
\dfrac{5}{15} &= \dfrac{5 \times 1}{5 \times 3}\\
              &= \dfrac{1}{3}
\end{align*}
\end{minipage}\hfill
\begin{minipage}[t]{0.31\textwidth}
\textbf{b)} 
\begin{align*}
\dfrac{20}{30} &= \dfrac{10 \times 2}{10 \times 3}\\
               &= \dfrac{2}{3}
\end{align*}
\end{minipage}\hfill
\begin{minipage}[t]{0.31\textwidth}
\textbf{c)} 
\begin{align*}
\dfrac{105}{165} &= \dfrac{15 \times 7}{15 \times 11}\\
                 &= \dfrac{7}{11}
\end{align*}
\end{minipage}

% =====================================================================
\vspace{1cm}\exercice
Le nombre de paquets doit diviser à la fois 150 et 105 : c'est un diviseur commun à
150 et 105. Pour en faire un maximum, on cherche le plus grand diviseur commun à 150 et 105.

Avec les décompositions en facteurs premiers :
$150 = 2 \times 3 \times 5 \times 5$ et $105 = 3 \times 5 \times 7$.
Le plus grand diviseur commun à 150 et 105 est le produit des facteurs premiers communs :
3 \times 5
                      = 15

Hulk et Batman peuvent donc faire au maximum \textbf{15 paquets},
chacun contenant  10 cookies et 7 tartelettes.

% =====================================================================
\vspace{1cm}\exercice
On utilise les critères de divisibilité :
\newline \textbf{par 2} : le chiffre des unités est 0, 2, 4, 6 ou 8 ;
\newline \textbf{par 3} : la somme des chiffres est divisible par 3
(sommes des chiffres : $28$ pour 594\,370 ; $36$ pour 444\,444\,444 ; $29$ pour 726\,815) ;
\newline \textbf{par 5} : le chiffre des unités est 0 ou 5 ;
\newline \textbf{par 10} : le chiffre des unités est 0.

\renewcommand{\arraystretch}{1.4}
\setlength{\tabcolsep}{2pt}
\begin{center}
\small
\begin{tabular}{|c|c|c|c|c|}
\hline
 & \textbf{Divisible par 2} & \textbf{Divisible par 3} & \textbf{Divisible par 5} & \textbf{Divisible par 10} \\
\hline
594\,370 & Oui & Non & Oui & Oui \\
\hline
444\,444\,444 & Oui & Oui & Non & Non \\
\hline
726\,815 & Non & Non & Oui & Non \\
\hline
\end{tabular}
\end{center}

\newpage

% =====================================================================
% VERSION 5
% =====================================================================
\setcounter{exercice}{0}

% ----- En-tête -----
\begin{center}
  \underline{\Large\textbf{Correction du contrôle séquence 13 version 5}}\\[0.3em]
 
\end{center}

% =====================================================================
\vspace{3cm}\exercice
\begin{enumerate}
  \item \textbf{Oui.} $45  = 5 \times 9$ 
  \newline 5 est donc un diviseur de 45.
  \item \textbf{Non.} $187 \div 6 \approx 31{,}17$
  \newline le quotient n'est pas un entier,
        donc 187 n'est pas un multiple de 6.
  \item \textbf{Non.}  $25 < 75$, donc 25 n'est pas un multiple de 75.
  \item \textbf{Non.} $54 \div 8 = 6{,}75$ 
  \newline le quotient n'est pas un entier,
        donc 8 n'est pas un diviseur de 54.
\end{enumerate}

% =====================================================================
\vspace{1cm}\exercice
\begin{enumerate}
  \item \textbf{2 et 3} (acceptés aussi : 5 ; 7 ; 11 ; \dots).
        Un nombre premier a exactement deux diviseurs : 1 et lui-même.
  \item \textbf{4 et 6} (acceptés aussi : 1 ; 8 ; 9 ; \dots).
      
  \item Dans 14 ; 22 ; 11 ; 27 ; 49 ; 58 ; 83 ; 121, les deux nombres premiers sont
        \textbf{11} et \textbf{83}.
       
\end{enumerate}

% =====================================================================
\vspace{1cm}\exercice
\textbf{a) } $21 = 3 \times 7$, 
donc les diviseurs de 21 sont $\boxed{1\ ;\ 21\ ;\ 3\ ;\ 7\ }$.

\textbf{b) } $28 = 2 \times 2 \times 7$,
donc les diviseurs de 28 sont
$\boxed{1\ ;\ 28\ ;\ 2\ ;\ 7\ ;\ 4\ ;\ 14\ }$.

% =====================================================================
\vspace{1cm}\exercice


\noindent\begin{minipage}[t]{0.48\textwidth}
\textbf{a)}
\begin{align*}
26 &= 2 \times 13
\end{align*}

\textbf{b)}
\begin{align*}
63 &= 3 \times 21\\
   &= 3 \times 3 \times 7
\end{align*}
\end{minipage}\hfill
\begin{minipage}[t]{0.48\textwidth}
\textbf{c)}
\begin{align*}
140 &= 2 \times 70\\
    &= 2 \times 2 \times 35\\
    &= 2 \times 2 \times 5 \times 7\\
\end{align*}

\textbf{d)}
\begin{align*}
182 &= 2 \times 91\\
    &= 2 \times 7 \times 13
\end{align*}
\end{minipage}

% =====================================================================
\exercice
\noindent\begin{minipage}[t]{0.31\textwidth}
\textbf{a)} 
\begin{align*}
\dfrac{6}{18} &= \dfrac{6 \times 1}{6 \times 3}\\
              &= \dfrac{1}{3}
\end{align*}
\end{minipage}\hfill
\begin{minipage}[t]{0.31\textwidth}
\textbf{b)} 
\begin{align*}
\dfrac{10}{15} &= \dfrac{5 \times 2}{5 \times 3}\\
               &= \dfrac{2}{3}
\end{align*}
\end{minipage}\hfill
\begin{minipage}[t]{0.31\textwidth}
\textbf{c)} 
\begin{align*}
\dfrac{75}{165} &= \dfrac{15 \times 5}{15 \times 11}\\
                &= \dfrac{5}{11}
\end{align*}
\end{minipage}

% =====================================================================
\vspace{1cm}\exercice
Le nombre de paquets doit diviser à la fois 135 et 105 : c'est un diviseur commun à
135 et 105. Pour en faire un maximum, on cherche le plus grand diviseur commun à 135 et 105.

Avec les décompositions en facteurs premiers :
$135 = 3 \times 3 \times 3 \times 5$ et $105 = 3 \times 5 \times 7$.
Le plus grand diviseur commun à 135 et 105 est le produit des facteurs premiers communs :
3 \times 5
                      = 15

Hulk et Batman peuvent donc faire au maximum \textbf{15 paquets},
chacun contenant  9 cookies et 7 tartelettes.

% =====================================================================
\vspace{1cm}\exercice
On utilise les critères de divisibilité :
\newline \textbf{par 2} : le chiffre des unités est 0, 2, 4, 6 ou 8 ;
\newline \textbf{par 3} : la somme des chiffres est divisible par 3
(sommes des chiffres : $31$ pour 671\,890 ; $54$ pour 666\,666\,666 ; $40$ pour 675\,985) ;
\newline \textbf{par 5} : le chiffre des unités est 0 ou 5 ;
\newline \textbf{par 10} : le chiffre des unités est 0.

\renewcommand{\arraystretch}{1.4}
\setlength{\tabcolsep}{2pt}
\begin{center}
\small
\begin{tabular}{|c|c|c|c|c|}
\hline
 & \textbf{Divisible par 2} & \textbf{Divisible par 3} & \textbf{Divisible par 5} & \textbf{Divisible par 10} \\
\hline
671\,890 & Oui & Non & Oui & Oui \\
\hline
666\,666\,666 & Oui & Oui & Non & Non \\
\hline
675\,985 & Non & Non & Oui & Non \\
\hline
\end{tabular}
\end{center}

\newpage

% =====================================================================
% VERSION 6
% =====================================================================
\setcounter{exercice}{0}

% ----- En-tête -----
\begin{center}
  \underline{\Large\textbf{Correction du contrôle séquence 13 version 6}}\\[0.3em]
 
\end{center}

% =====================================================================
\vspace{3cm}\exercice
\begin{enumerate}
  \item \textbf{Oui.} $32  = 4 \times 8$ 
  \newline 4 est donc un diviseur de 32.
  \item \textbf{Non.} $153 \div 4 = 38{,}25$
  \newline le quotient n'est pas un entier,
        donc 153 n'est pas un multiple de 4.
  \item \textbf{Non.}  $18 < 54$, donc 18 n'est pas un multiple de 54.
  \item \textbf{Non.} $60 \div 7 \approx 8{,}57$
  \newline le quotient n'est pas un entier,
        donc 7 n'est pas un diviseur de 60.
\end{enumerate}

% =====================================================================
\vspace{1cm}\exercice
\begin{enumerate}
  \item \textbf{2 et 3} (acceptés aussi : 5 ; 7 ; 11 ; \dots).
        Un nombre premier a exactement deux diviseurs : 1 et lui-même.
  \item \textbf{4 et 6} (acceptés aussi : 1 ; 8 ; 9 ; \dots).
      
  \item Dans 15 ; 26 ; 7 ; 21 ; 51 ; 66 ; 73 ; 95, les deux nombres premiers sont
        \textbf{7} et \textbf{73}.
       
\end{enumerate}

% =====================================================================
\vspace{1cm}\exercice
\textbf{a) } $22 = 2 \times 11$, 
donc les diviseurs de 22 sont $\boxed{1\ ;\ 22\ ;\ 2\ ;\ 11\ }$.

\textbf{b) } $63 = 3 \times 3 \times 7$,
donc les diviseurs de 63 sont
$\boxed{1\ ;\ 63\ ;\ 3\ ;\ 7\ ;\ 9\ ;\ 21\ }$.

% =====================================================================
\vspace{1cm}\exercice


\noindent\begin{minipage}[t]{0.48\textwidth}
\textbf{a)}
\begin{align*}
33 &= 3 \times 11
\end{align*}

\textbf{b)}
\begin{align*}
45 &= 3 \times 15\\
   &= 3 \times 3 \times 5
\end{align*}
\end{minipage}\hfill
\begin{minipage}[t]{0.48\textwidth}
\textbf{c)}
\begin{align*}
220 &= 2 \times 110\\
    &= 2 \times 2 \times 55\\
    &= 2 \times 2 \times 5 \times 11\\
\end{align*}

\textbf{d)}
\begin{align*}
273 &= 3 \times 91\\
    &= 3 \times 7 \times 13
\end{align*}
\end{minipage}

% =====================================================================
\exercice
\noindent\begin{minipage}[t]{0.31\textwidth}
\textbf{a)} 
\begin{align*}
\dfrac{4}{12} &= \dfrac{4 \times 1}{4 \times 3}\\
              &= \dfrac{1}{3}
\end{align*}
\end{minipage}\hfill
\begin{minipage}[t]{0.31\textwidth}
\textbf{b)} 
\begin{align*}
\dfrac{18}{27} &= \dfrac{9 \times 2}{9 \times 3}\\
               &= \dfrac{2}{3}
\end{align*}
\end{minipage}\hfill
\begin{minipage}[t]{0.31\textwidth}
\textbf{c)} 
\begin{align*}
\dfrac{30}{195} &= \dfrac{15 \times 2}{15 \times 13}\\
                &= \dfrac{2}{13}
\end{align*}
\end{minipage}

% =====================================================================
\vspace{1cm}\exercice
Le nombre de paquets doit diviser à la fois 165 et 120 : c'est un diviseur commun à
165 et 120. Pour en faire un maximum, on cherche le plus grand diviseur commun à 165 et 120.

Avec les décompositions en facteurs premiers :
$165 = 3 \times 5 \times 11$ et $120 = 2 \times 2 \times 2 \times 3 \times 5$.
Le plus grand diviseur commun à 165 et 120 est le produit des facteurs premiers communs :
3 \times 5
                      = 15

Hulk et Batman peuvent donc faire au maximum \textbf{15 paquets},
chacun contenant  11 cookies et 8 tartelettes.

% =====================================================================
\vspace{1cm}\exercice
On utilise les critères de divisibilité :
\newline \textbf{par 2} : le chiffre des unités est 0, 2, 4, 6 ou 8 ;
\newline \textbf{par 3} : la somme des chiffres est divisible par 3
(sommes des chiffres : $32$ pour 728\,960 ; $63$ pour 777\,777\,777 ; $26$ pour 341\,765) ;
\newline \textbf{par 5} : le chiffre des unités est 0 ou 5 ;
\newline \textbf{par 10} : le chiffre des unités est 0.

\renewcommand{\arraystretch}{1.4}
\setlength{\tabcolsep}{2pt}
\begin{center}
\small
\begin{tabular}{|c|c|c|c|c|}
\hline
 & \textbf{Divisible par 2} & \textbf{Divisible par 3} & \textbf{Divisible par 5} & \textbf{Divisible par 10} \\
\hline
728\,960 & Oui & Non & Oui & Oui \\
\hline
777\,777\,777 & Non & Oui & Non & Non \\
\hline
341\,765 & Non & Non & Oui & Non \\
\hline
\end{tabular}
\end{center}

\newpage

% =====================================================================
% VERSION 7
% =====================================================================
\setcounter{exercice}{0}

% ----- En-tête -----
\begin{center}
  \underline{\Large\textbf{Correction du contrôle séquence 13 version 7}}\\[0.3em]
 
\end{center}

% =====================================================================
\vspace{3cm}\exercice
\begin{enumerate}
  \item \textbf{Oui.} $48  = 6 \times 8$ 
  \newline 6 est donc un diviseur de 48.
  \item \textbf{Non.} $184 \div 9 \approx 20{,}44$
  \newline le quotient n'est pas un entier,
        donc 184 n'est pas un multiple de 9.
  \item \textbf{Non.}  $21 < 63$, donc 21 n'est pas un multiple de 63.
  \item \textbf{Non.} $46 \div 4 = 11{,}5$ 
  \newline le quotient n'est pas un entier,
        donc 4 n'est pas un diviseur de 46.
\end{enumerate}

% =====================================================================
\vspace{1cm}\exercice
\begin{enumerate}
  \item \textbf{2 et 3} (acceptés aussi : 5 ; 7 ; 11 ; \dots).
        Un nombre premier a exactement deux diviseurs : 1 et lui-même.
  \item \textbf{4 et 6} (acceptés aussi : 1 ; 8 ; 9 ; \dots).
      
  \item Dans 16 ; 25 ; 5 ; 33 ; 65 ; 76 ; 89 ; 87, les deux nombres premiers sont
        \textbf{5} et \textbf{89}.
       
\end{enumerate}

% =====================================================================
\vspace{1cm}\exercice
\textbf{a) } $26 = 2 \times 13$, 
donc les diviseurs de 26 sont $\boxed{1\ ;\ 26\ ;\ 2\ ;\ 13\ }$.

\textbf{b) } $18 = 2 \times 3 \times 3$,
donc les diviseurs de 18 sont
$\boxed{1\ ;\ 18\ ;\ 2\ ;\ 3\ ;\ 6\ ;\ 9\ }$.

% =====================================================================
\vspace{1cm}\exercice


\noindent\begin{minipage}[t]{0.48\textwidth}
\textbf{a)}
\begin{align*}
35 &= 5 \times 7
\end{align*}

\textbf{b)}
\begin{align*}
28 &= 2 \times 14\\
   &= 2 \times 2 \times 7
\end{align*}
\end{minipage}\hfill
\begin{minipage}[t]{0.48\textwidth}
\textbf{c)}
\begin{align*}
260 &= 2 \times 130\\
    &= 2 \times 2 \times 65\\
    &= 2 \times 2 \times 5 \times 13\\
\end{align*}

\textbf{d)}
\begin{align*}
286 &= 2 \times 143\\
    &= 2 \times 11 \times 13
\end{align*}
\end{minipage}

% =====================================================================
\exercice
\noindent\begin{minipage}[t]{0.31\textwidth}
\textbf{a)} 
\begin{align*}
\dfrac{6}{21} &= \dfrac{3 \times 2}{3 \times 7}\\
              &= \dfrac{2}{7}
\end{align*}
\end{minipage}\hfill
\begin{minipage}[t]{0.31\textwidth}
\textbf{b)} 
\begin{align*}
\dfrac{24}{36} &= \dfrac{12 \times 2}{12 \times 3}\\
               &= \dfrac{2}{3}
\end{align*}
\end{minipage}\hfill
\begin{minipage}[t]{0.31\textwidth}
\textbf{c)} 
\begin{align*}
\dfrac{75}{195} &= \dfrac{15 \times 5}{15 \times 13}\\
                &= \dfrac{5}{13}
\end{align*}
\end{minipage}

% =====================================================================
\vspace{1cm}\exercice
Le nombre de paquets doit diviser à la fois 180 et 105 : c'est un diviseur commun à
180 et 105. Pour en faire un maximum, on cherche le plus grand diviseur commun à 180 et 105.

Avec les décompositions en facteurs premiers :
$180 = 2 \times 2 \times 3 \times 3 \times 5$ et $105 = 3 \times 5 \times 7$.
Le plus grand diviseur commun à 180 et 105 est le produit des facteurs premiers communs :
3 \times 5
                      = 15

Hulk et Batman peuvent donc faire au maximum \textbf{15 paquets},
chacun contenant  12 cookies et 7 tartelettes.

% =====================================================================
\vspace{1cm}\exercice
On utilise les critères de divisibilité :
\newline \textbf{par 2} : le chiffre des unités est 0, 2, 4, 6 ou 8 ;
\newline \textbf{par 3} : la somme des chiffres est divisible par 3
(sommes des chiffres : $29$ pour 845\,930 ; $72$ pour 888\,888\,888 ; $35$ pour 953\,585) ;
\newline \textbf{par 5} : le chiffre des unités est 0 ou 5 ;
\newline \textbf{par 10} : le chiffre des unités est 0.

\renewcommand{\arraystretch}{1.4}
\setlength{\tabcolsep}{2pt}
\begin{center}
\small
\begin{tabular}{|c|c|c|c|c|}
\hline
 & \textbf{Divisible par 2} & \textbf{Divisible par 3} & \textbf{Divisible par 5} & \textbf{Divisible par 10} \\
\hline
845\,930 & Oui & Non & Oui & Oui \\
\hline
888\,888\,888 & Oui & Oui & Non & Non \\
\hline
953\,585 & Non & Non & Oui & Non \\
\hline
\end{tabular}
\end{center}

\newpage

% =====================================================================
% VERSION 8
% =====================================================================
\setcounter{exercice}{0}

% ----- En-tête -----
\begin{center}
  \underline{\Large\textbf{Correction du contrôle séquence 13 version 8}}\\[0.3em]
 
\end{center}

% =====================================================================
\vspace{3cm}\exercice
\begin{enumerate}
  \item \textbf{Oui.} $64  = 8 \times 8$ 
  \newline 8 est donc un diviseur de 64.
  \item \textbf{Non.} $222 \div 7 \approx 31{,}71$
  \newline le quotient n'est pas un entier,
        donc 222 n'est pas un multiple de 7.
  \item \textbf{Non.}  $35 < 105$, donc 35 n'est pas un multiple de 105.
  \item \textbf{Non.} $56 \div 9 \approx 6{,}22$
  \newline le quotient n'est pas un entier,
        donc 9 n'est pas un diviseur de 56.
\end{enumerate}

% =====================================================================
\vspace{1cm}\exercice
\begin{enumerate}
  \item \textbf{2 et 3} (acceptés aussi : 5 ; 7 ; 11 ; \dots).
        Un nombre premier a exactement deux diviseurs : 1 et lui-même.
  \item \textbf{4 et 6} (acceptés aussi : 1 ; 8 ; 9 ; \dots).
      
  \item Dans 18 ; 27 ; 2 ; 35 ; 77 ; 82 ; 67 ; 91, les deux nombres premiers sont
        \textbf{2} et \textbf{67}.
       
\end{enumerate}

% =====================================================================
\vspace{1cm}\exercice
\textbf{a) } $33 = 3 \times 11$, 
donc les diviseurs de 33 sont $\boxed{1\ ;\ 33\ ;\ 3\ ;\ 11\ }$.

\textbf{b) } $75 = 3 \times 5 \times 5$,
donc les diviseurs de 75 sont
$\boxed{1\ ;\ 75\ ;\ 3\ ;\ 5\ ;\ 15\ ;\ 25\ }$.

% =====================================================================
\vspace{1cm}\exercice


\noindent\begin{minipage}[t]{0.48\textwidth}
\textbf{a)}
\begin{align*}
39 &= 3 \times 13
\end{align*}

\textbf{b)}
\begin{align*}
98 &= 2 \times 49\\
   &= 2 \times 7 \times 7
\end{align*}
\end{minipage}\hfill
\begin{minipage}[t]{0.48\textwidth}
\textbf{c)}
\begin{align*}
308 &= 2 \times 154\\
    &= 2 \times 2 \times 77\\
    &= 2 \times 2 \times 7 \times 11\\
\end{align*}

\textbf{d)}
\begin{align*}
385 &= 5 \times 77\\
    &= 5 \times 7 \times 11
\end{align*}
\end{minipage}

% =====================================================================
\exercice
\noindent\begin{minipage}[t]{0.31\textwidth}
\textbf{a)} 
\begin{align*}
\dfrac{9}{12} &= \dfrac{3 \times 3}{3 \times 4}\\
              &= \dfrac{3}{4}
\end{align*}
\end{minipage}\hfill
\begin{minipage}[t]{0.31\textwidth}
\textbf{b)} 
\begin{align*}
\dfrac{28}{42} &= \dfrac{14 \times 2}{14 \times 3}\\
               &= \dfrac{2}{3}
\end{align*}
\end{minipage}\hfill
\begin{minipage}[t]{0.31\textwidth}
\textbf{c)} 
\begin{align*}
\dfrac{45}{105} &= \dfrac{15 \times 3}{15 \times 7}\\
                &= \dfrac{3}{7}
\end{align*}
\end{minipage}

% =====================================================================
\vspace{1cm}\exercice
Le nombre de paquets doit diviser à la fois 105 et 150 : c'est un diviseur commun à
105 et 150. Pour en faire un maximum, on cherche le plus grand diviseur commun à 105 et 150.

Avec les décompositions en facteurs premiers :
$105 = 3 \times 5 \times 7$ et $150 = 2 \times 3 \times 5 \times 5$.
Le plus grand diviseur commun à 105 et 150 est le produit des facteurs premiers communs :
3 \times 5
                      = 15

Hulk et Batman peuvent donc faire au maximum \textbf{15 paquets},
chacun contenant  7 cookies et 10 tartelettes.

% =====================================================================
\vspace{1cm}\exercice
On utilise les critères de divisibilité :
\newline \textbf{par 2} : le chiffre des unités est 0, 2, 4, 6 ou 8 ;
\newline \textbf{par 3} : la somme des chiffres est divisible par 3
(sommes des chiffres : $29$ pour 917\,480 ; $81$ pour 999\,999\,999 ; $31$ pour 459\,175) ;
\newline \textbf{par 5} : le chiffre des unités est 0 ou 5 ;
\newline \textbf{par 10} : le chiffre des unités est 0.

\renewcommand{\arraystretch}{1.4}
\setlength{\tabcolsep}{2pt}
\begin{center}
\small
\begin{tabular}{|c|c|c|c|c|}
\hline
 & \textbf{Divisible par 2} & \textbf{Divisible par 3} & \textbf{Divisible par 5} & \textbf{Divisible par 10} \\
\hline
917\,480 & Oui & Non & Oui & Oui \\
\hline
999\,999\,999 & Non & Oui & Non & Non \\
\hline
459\,175 & Non & Non & Oui & Non \\
\hline
\end{tabular}
\end{center}
\end{document}
